\begin{proof}\textbf{of Theorem~\ref{thm: FTP MSE RKHS}.}
Let $\bar{\mathbf f}=n^{-1}\sum_{i=1}^n\mathbf f_i$. Since $K_T$ is positive definite, $K_T^{-1}$ and $K_T^{1/2}$ are well-defined.\\
\textcolor{red}{a.} 
We first show that for any $f_i\in \mathcal{H}_K$ with $\|f_i\|_{\mathcal{H}_K}\leq B$, we can always find another function $g_i$ in the span of $\{K(t_1,t),\ldots,K(t_T,t)\}$ such that 
\[
f_i(t_j)=g_i(t_j),\ \|g_i\|_{\mathcal{H}_K}\leq B.
\]
For each $i$, consider the minimum-norm interpolation problem
\[
g_i
=
\argmin_{g\in\mathcal H_K}
\left\{
\|g\|_{\mathcal H_K}:
g(t_j)=f_i(t_j),\ j=1,\ldots,T
\right\}.
\]
By the representer theorem,
\[
g_i(t)
=
\sum_{j=1}^T
(\gamma_i)_jK(t_j,t),
\quad
\gamma_i=K_T^{-1}\mathbf f_i.
\]
Hence
\[
\|g_i\|_{\mathcal H_K}^2
=
\gamma_i^\top K_T\gamma_i
=
\mathbf f_i^\top K_T^{-1}\mathbf f_i.
\]
Since $f_i$ itself is feasible for the interpolation problem,
\[
\|g_i\|_{\mathcal H_K}
\le
\|f_i\|_{\mathcal H_K}
\le B.
\]
Therefore,
\begin{equation}
\mathbf f_i^\top K_T^{-1}\mathbf f_i
\le B^2.
\label{eq: RKHS_min_norm_true_bound}
\end{equation}
Define
\[
\mathbf \beta_i
:=
K_T^{-1/2}\mathbf f_i.
\]
Then Equation~\eqref{eq: RKHS_min_norm_true_bound} gives
\begin{equation}
\|g_i\|_{\mathcal H_K} = \|\beta_i\|_2\le B.
\label{eq:true_transformed_norm}
\end{equation}
\textcolor{red}{b.} 
We then derive the form of the constrained and unconstrained solution to \begin{equation}
    \hat{f}_i^{(c)} 
    = \argmin_{\|f\|_{\mathcal{H}_K}\le B'} 
    \dfrac{1}{T}\sum_{j=1}^T [x_{ij}-f(t_j)]^2.
\label{eq: constrained_FTP_RKHS}
\end{equation}
and perform coefficient transformation for the convenience of following proof.\\ 
By the Representer theorem, the constrained estimator has the form
\[
\hat f_i^{(c)}(t)
=
\sum_{j=1}^T
(\hat\alpha_i^{(c)})_jK(t_j,t),
\]
and its RKHS norm satisfies
\[
\|\hat f_i^{(c)}\|_{\mathcal H_K}^2
=
\hat\alpha_i^{(c)\top}
K_T
\hat\alpha_i^{(c)}.
\]
Consequently, Equation~\eqref{eq: constrained_FTP_RKHS} is equivalent to
\begin{equation}
    \hat\alpha_i^{(c)}
    =
    \argmin_{\alpha^\top K_T\alpha\le {B'}^2}
    \|x_i-K_T\alpha\|_2^2.
    \label{eq: constrained_alpha_RKHS}
\end{equation}
Consider the corresponding unconstrained estimator
\[
\hat f_i^{(u)}
=
\argmin_{f\in\mathcal H_K}
\frac1T\sum_{j=1}^T
[x_{ij}-f(t_j)]^2.
\]
The minimum-norm unconstrained least-squares solution has the form
\begin{equation}
    \begin{aligned}
        \hat f_i^{(u)}(t)
        &=
        \sum_{j=1}^T [\hat\alpha_i^{(u)}]_jK(t_j,t),\\
        \hat\alpha_i^{(u)}
        &=
        K_T^{-1}x_i.
    \end{aligned}
\end{equation}
Indeed, the fitted-value vector is
\[
\hat{\mathbf{f}}_i^{(u)}
=
K_T\hat\alpha_i^{(u)}
=
x_i.
\]
Define
\[
\hat{\beta}_i^{(u)}
:=
K_T^{1/2}\hat\alpha_i^{(u)}
=
K_T^{-1/2}x_i
\text{\quad and\quad}
\hat{\beta} := K_T^{1/2}\alpha.
\]
We have
\[
\alpha^\top K_T\alpha\le {B'}^2
\Longleftrightarrow
\|\beta\|_2 = \|K_T^{1/2}\alpha\|_2\le B'.
\]
Hence, we can rewrite the optimization problem in Equation~\eqref{eq: constrained_alpha_RKHS} as
\begin{equation}
    \hat\beta_i^{(c)}
    =
    \argmin_{\|\beta\|_2\le B'}
    \|x_i-K_T^{1/2}\beta\|_2^2,
    \quad \hat\beta_i^{(c)} = K_T^{1/2}\hat\alpha_i^{(c)}.
    \label{eq: constrained_beta_RKHS}
\end{equation}
It follows
\[
\hat f_i^{(c)}(t)
=
\sum_{j=1}^T
(\hat\alpha_i^{(c)})_jK(t_j,t)
=
\sum_{j=1}^T
(K_T^{-1/2}\hat\beta_i^{(c)})_jK(t_j,t).
\]
\textcolor{red}{c.} 
Next, we introduce the metric projection and follow similar steps in the proof of Theorem~\ref{thm: FTP IMSE}.\\
Denote
\[
\mathcal C_{B'}
:=
\left\{
\beta\in\mathbb R^T: \|\beta\|_2\leq B'
\right\}.
\]
By substituting $\Phi = K_T^{1/2}$, $M = \Phi^\top \Phi = \Phi^2 = K_T$, we have
\begin{equation}
\hat{\beta}_i^{(c)}
=
\argmin_{\beta\in\mathcal C_{B'}}
(\beta-\hat{\beta}_i^{(u)})^\top
M
(\beta-\hat{\beta}_i^{(u)})
:= P_{\mathcal C_{B'}}^{M}
\left(
\hat{\beta}_i^{(u)}
\right).
\label{eq:constrained_projection_RKHS}
\end{equation}
where $P_{\mathcal C_{B'}}^{M}$ denotes the metric projection onto
$\mathcal C_{B'}$ with respect to the $M$-norm.\\
Since the projection is nonexpansive, and $\|\beta_i\|_2\leq B\leq B'$, we have
\begin{equation}
    \begin{aligned}
        \left\|
        \hat{\beta}_i^{(c)}-\beta_i
        \right\|_M
        =
        \left\|
        P_{\mathcal C_{B'}}^{M}
        \left(
        \hat{\beta}_i^{(u)}
        \right)
        -
        P_{\mathcal C_{B'}}^{M}(\beta_i)
        \right\|_M
        \leq
        \left\|
        \hat{\beta}_i^{(u)}-\beta_i
        \right\|_M.
    \end{aligned}
    \label{eq: M_norm_nonexpansive_RKHS}
\end{equation}
\textcolor{red}{d.} 
We denote convenient notations as follows
\begin{equation}
    \begin{aligned}
        e_i = \hat{\beta}_i^{(c)} - \beta_i, &&
        \xi_i = \hat{\beta}_i^{(u)} - \beta_i, &&
        r_i = \hat{\beta}_i^{(c)} - \hat{\beta}_i^{(u)}, &&
        b_i = \E e_i.
    \end{aligned}
    \label{eq: notations_RKHS}
\end{equation}
Then Equation~\eqref{eq: M_norm_nonexpansive_RKHS} becomes
\begin{equation}
    \|e_i\|_M\leq \|\xi_i\|_M.
    \label{eq: M_norm_nonexpansive_RKHS_simplified}
\end{equation}
Since
\[
x_i=\mathbf f_i+\epsilon_i,\quad
\epsilon_i\sim N(0,\sigma_0^2 I_T),
\]
we have
\[
\beta_i^{(u)}
=K_T^{-1/2} x_i
=
\Phi^{-1}x_i
=
\beta_i+\xi_i,
\]
where
\[
\xi_i
=
\Phi^{-1}\epsilon_i,\quad
\xi_i
\sim
N(0,\sigma_0^2M^{-1}).
\]
Consequently,
\begin{align}
\mathbb E
\left\|
\xi_i
\right\|_2^2
=
\operatorname{tr}
\left[
\operatorname{Cov}
\left(
\xi_i
\right)
\right]
=
\sigma_0^2\operatorname{tr}(M^{-1}),
\quad
\mathbb E
\left\|
\epsilon_i
\right\|_2^2
=\sigma_0^2 T.
\label{eq: unconstrained_mse_RKHS}
\end{align}
Then we can expand the MSE as
\begin{align}
\mathrm{MSE}_{\mathrm{FTP}}
&= 
\dfrac{1}{T} \sum_{j=1}^T \mathbb{E} 
[(\widetilde{f}_{\mathrm{FTP}}(t_j)-\bar{f}(t_j))^2]\nonumber\\
&= 
\dfrac{1}{T} \mathbb{E} 
\|\widetilde{\mathbf{f}}_{\mathrm{FTP}} - \bar{\mathbf{f}}\|_2^2\nonumber\\
&= 
\dfrac{1}{T} \mathbb{E} 
\|\overline{\hat{\mathbf{f}}^{(c)}} - \bar{\mathbf{f}} + \sigma \mathbf{Z}\|_2^2
\label{eq: direct expand MSE}.
\end{align}
Notice that $Z$ is a mean zero Gaussian process with kernel $K$, we have
\[
\mathbf{Z} = (Z(t_1),\ldots,Z(t_T))^\top
\sim \mathcal{N}(0,M),
\]
\[
\mathbb E\|\mathbf{Z}\|_2^2 = \operatorname{tr}(M).
\]
Thus, we have
\begin{align}
\mathrm{MSE}_{\mathrm{FTP}}
&= 
\dfrac{1}{T} \mathbb{E} 
\|\overline{\hat{\mathbf{f}}^{(c)}} - \bar{\mathbf{f}} + \sigma \mathbf{Z}\|_2^2\nonumber\\
&= 
\dfrac{1}{T} 
\Bra{\mathbb{E} 
\|\overline{\hat{\mathbf{f}}^{(c)}} - \bar{\mathbf{f}}\|_2^2 
+ 
\mathbb{E} \|\sigma \mathbf{Z}\|_2^2 + 0}
\nonumber\\
&= 
\dfrac{1}{T} 
\mathbb{E} 
\|\overline{\hat{\mathbf{f}}^{(c)}} - \bar{\mathbf{f}}\|_2^2 
+ 
\dfrac{4{B'}^2}{n^2\mu^2 T} \operatorname{tr}(M).
\label{eq: expanded MSE}
\end{align}
Note that $e_i$'s are independent, we can further expand the first term as
\begin{equation}
    \begin{aligned}
        \mathbb{E} 
        \|\overline{\hat{\mathbf{f}}^{(c)}} - \bar{\mathbf{f}}\|_2^2 
        &=
        \mathbb{E} 
        \|\overline{\hat{\mathbf{f}}^{(c)}} - \bar{\mathbf{g}}\|_2^2\\
        &=
        \mathbb{E} 
        \|\Phi \overline{\hat{\beta}^{(c)}} - \Phi\bar{\beta}\|_2^2\\
        &=
        \mathbb{E} 
        \|\overline{\hat{\beta}^{(c)}} - \bar{\beta}\|_{M}^2\\
        &=
        \dfrac{1}{n^2} \mathbb E \Bra{
        \sum_{i=1}^n \|e_i\|_M^2 + \sum_{i\ne j} e_i^\top  M e_j}\\
        &=
        \dfrac{1}{n^2} \Bra{
        \sum_{i=1}^n \mathbb E \|e_i\|_M^2
        + \sum_{i\ne j} \mathbb E e_i^\top M \mathbb E e_j
        }.
    \end{aligned}
    \label{eq: expand bias and clipping}
\end{equation}
Recall Equation~\eqref{eq: M_norm_nonexpansive_RKHS_simplified} and Equation~\eqref{eq: unconstrained_mse_RKHS}, we have
\begin{equation}
    \begin{aligned}
        \dfrac{1}{n^2} \sum_{i=1}^n \mathbb E \|e_i\|_M^2 
        \leq
        \dfrac{1}{n^2} \sum_{i=1}^n \mathbb E \|\xi_i\|_M^2
        =
        \dfrac{1}{n^2} \sum_{i=1}^n \mathbb E \|\epsilon_i\|_2^2
        =
        \dfrac{\sigma_0^2T}{n}.
    \end{aligned}
\label{eq: noise term_RKHS}
\end{equation}
\textcolor{red}{e.} 
Finally, we analyze the cross-term.\\
Recall the definition of $\hat{\beta}_i^{(c)}$ given by Equation~\eqref{eq:constrained_projection_RKHS}, we have
\begin{equation}
\|r_i\|_M
=
\left\|
\hat{\beta}_i^{(c)} - \hat{\beta}_i^{(u)}
\right\|_M
\le 
\left\|
\beta_i - \hat{\beta}_i^{(u)}
\right\|_M
=
\|\xi_i\|_M.
\end{equation}
Since $\hat{\beta}_i^{(u)}$ is the unconstrained minimizer, we know that
$r_i=0$ if $\|\hat{\beta}_i^{(u)}\|_2\le B'$.\\
Therefore, we can rewrite
\[
r_i 
= \Bra{
\hat{\beta}_i^{(c)} - \hat{\beta}_i^{(u)}
}\cdot \mathbf{1} \Set{\|\hat{\beta}_i^{(u)}\|_2>B'}.
\]
It follows that
\begin{align*}
\|r_i\|_M^2
=&\|\hat{\beta}_i^{(c)} - \hat{\beta}_i^{(u)}\|_M^2 \cdot \mathbf{1} \Set{\|\hat{\beta}_i^{(u)}\|_2>B'}\\
\le&
\|\xi_i\|_M^2
\cdot \mathbf{1} \Set{\|\hat{\beta}_i^{(u)}\|_2>B'}\\
=& \|\epsilon_i\|_2^2
\cdot \mathbf{1} \Set{\|\hat{\beta}_i^{(u)}\|_2>B'}.
\end{align*}
Thus, we can obtain that
\begin{equation}
    \|b_i\|_M^2
    =
    \| \mathbb E e_i \|_M^2
    =
    \| \mathbb E r_i \|_M^2
    \le
    \mathbb E \|r_i\|_M^2
    \le
    \mathbb E \SqBra{
    \|\epsilon_i\|_2^2
    \cdot \mathbf{1} \Set{\|\hat{\beta}_i^{(u)}\|_2>B'}
    }.
\label{eq: clipping bias bound_RKHS}
\end{equation}
By the Cauchy--Schwarz inequality, we have
\begin{equation}
    \mathbb E \SqBra{
    \|\epsilon_i\|_2^2
    \cdot \mathbf{1} \Set{\|\hat{\beta}_i^{(u)}\|_2>B'}
    }
    \leq
    \sqrt{\mathbb E \|\epsilon_i\|_2^4} \sqrt{\operatorname{Pr}(\|\hat{\beta}_i^{(u)}\|_2>B')}.
    \label{eq: Cauchy--Schwarz inequality_RKHS}
\end{equation}
Since $\epsilon_i\overset{i.i.d.}{\sim} N(0,\Sigma)$ with $\Sigma = \sigma_0^2 I_T$, we have
\begin{equation}
\sqrt{\mathbb E \|\epsilon_i\|_2^4}
=
\sqrt{(\operatorname{tr}\Sigma)^2 + 2\operatorname{tr}(\Sigma^{2})}
= \sigma_0^2 \sqrt{T^2+2T}.
\label{eq: 4th moment guassian bound_RKHS}
\end{equation}
Recall $\|\beta_i\|_2\le B$, denote $\delta=B'-B$, then 
$\Set{\|\hat{\beta}_i^{(u)}\|_2>B'}$ implies 
$\Set{\|\hat{\beta}_i^{(u)}-\beta_i\|_2>\delta}$, thus
\begin{align*}
\operatorname{Pr}(\|\hat{\beta}_i^{(u)}\|_2>B')
&\le
\operatorname{Pr}(\|\hat{\beta}_i^{(u)}-\beta_i\|_2>\delta)
\\
&=\operatorname{Pr}(\|\xi_i\|_2>\delta).
\end{align*}
Let $\xi_i = \sigma_0 M^{-1/2}Z_i$, $Z_i\sim N(0, I_T)$, $\|Z_i\|_2^2\sim \chi^2_T$, we have
\[
\|\xi_i\|_2^2
= 
\sigma_0^2 Z_i^\top M^{-1} Z_i 
\le 
\sigma_0^2 \lambda_{\max}(M^{-1}) \|Z_i\|_2^2
=
\dfrac{\sigma_0^2}{\lambda_{\min}(M)} \|Z_i\|_2^2.
\]
Apply the $\chi^2$ tail bound from \citet{laurent2000adaptive}, we have
\begin{align}
\operatorname{Pr}(\|\hat{\beta}_i^{(u)}\|_2>B')
&\le
\operatorname{Pr}(\|\xi_i\|_2>\delta)
\nonumber\\
&\le
\operatorname{Pr}
\Bra{
\chi^2_T > \dfrac{\delta^2 \lambda_{\min}(M)}{\sigma_0^2}
}
\nonumber\\
&\le
\exp\SqBra{
-\dfrac{1}{2}\Bra{
\dfrac{\delta\sqrt{\lambda_{\min}(M)}}{\sigma_0}
- \sqrt{T}
}^2
}.
\label{eq: clipping probability bound_RKHS}
\end{align}
Combine Equation~\eqref{eq: clipping bias bound_RKHS},~\eqref{eq: Cauchy--Schwarz inequality_RKHS},~\eqref{eq: 4th moment guassian bound_RKHS} and~\eqref{eq: clipping probability bound_RKHS}, we obtain
\begin{equation}
    \|b_i\|_M^2
    \leq
    \sigma_0^2\sqrt{T^2+2T}\cdot
    \exp\SqBra{
    -\dfrac{1}{4}\Bra{
    \dfrac{\delta\sqrt{\lambda_{\min}(M)}}{\sigma_0}
    - \sqrt{T}
    }^2
    }.
    \label{eq: individual_bias_term_RKHS}
\end{equation}
By the Cauchy--Schwarz inequality, we have
\begin{equation}
\dfrac{1}{n^2} \sum_{i\ne j} \mathbb E e_i^\top M \mathbb E e_j
=
\dfrac{1}{n^2} \sum_{i\ne j} b_i^\top M b_j
\le
\Bra{1-\dfrac{1}{n}} \max_{i} \|b_i\|^2_M.
\label{eq: Cauchy--Schwarz_RKHS}
\end{equation}
Finally combine Equation~\eqref{eq: expanded MSE},~\eqref{eq: expand bias and clipping},~\eqref{eq: noise term_RKHS},~\eqref{eq: individual_bias_term_RKHS},~\eqref{eq: Cauchy--Schwarz_RKHS}, and recall $K_T=M$, we get
\begin{equation}
\begin{aligned}
\mathrm{MSE}_{\mathrm{FTP}}
\leq
&\frac{4{B'}^2\operatorname{tr}(K_T)}{n^2\mu^2T}
+
\dfrac{\sigma_0^2}{n}\\
&+
\Bra{1-\dfrac{1}{n}}
\sigma_0^2 \sqrt{\dfrac{T+2}{T}}
\exp\SqBra{
-\dfrac{1}{4}\Bra{
\dfrac{(B'-B)\sqrt{\lambda_{\min}(K_T)}}{\sigma_0}
- \sqrt{T}
}^2
}.
\end{aligned}
\label{eq: final_MSE_FTP bound}
\end{equation}
When $\sigma_0=0$ and $B'\ge B$, a feasible zero-loss fit exists, so every selected minimizer satisfies $\hat f_i^{(c)}(t_j)=f_i(t_j)$. Consequently, only privacy noise contributes to the grid MSE, giving $\operatorname{MSE}_{\mathrm{FTP}}=4(B')^2\operatorname{tr}(K_T)/(n^2\mu^2T)$.
\end{proof}
\begin{remark}\textup{\textbf{(Minimum-Norm Interpolator Problem)}}
    Although both the FTP and PTF procedures involve the constrained optimization
    problems in Equations~\eqref{eq: F_operator} and~\eqref{eq: weight least square},
    the role of the constraint is different in the noiseless setting. For the
    individual trajectories, when $\sigma_0=0$ we have
    $x_{ij}=f_i(t_j)$. Hence, the empirical least-squares loss can attain its
    global minimum value zero. In particular, the constrained problem
    \[
    \hat f_i
    =
    \argmin_{\substack{f\in\mathcal H_K\\
    \|f\|_{\mathcal H_K}\le B'}}
    \frac{1}{T}\sum_{j=1}^T
    \left[f(t_j)-x_{ij}\right]^2
    \]
    admits $f_i$ as a feasible zero-loss solution. Consequently, the constraint
    does not affect the optimal value of the individual fitting problem.\vspace{1em}\\
    A convenient representative of the zero-loss solutions is the
    minimum-norm interpolant
    \begin{equation}
    \hat f_i^{\mathrm{MN}}
    =
    \argmin_{f\in\mathcal H_K}
    \|f\|_{\mathcal H_K}
    \quad
    \text{subject to}
    \quad
    f(t_j)=x_{ij},\qquad j=1,\ldots,T.
    \label{eq: minimum norm interpolator}
    \end{equation}
    Let $K_T=\left[K(t_j,t_l)\right]_{j,l=1}^T$ be the kernel Gram matrix, which is positive definite for distinct sampling locations under the Gaussian kernel. A direct Hilbert-space projection argument shows that the minimum-norm interpolation belongs to the finite-dimensional span of the kernel sections:
    \[
    \hat f_i^{\mathrm{MN}}(t)
    =
    k_T(t)^\top K_T^{-1}x_i,
    \qquad
    k_T(t)
    =
    \left(K(t,t_1),\ldots,K(t,t_T)\right)^\top.
    \]
    Moreover, $\|\hat f_i^{\mathrm{MN}}\|_{\mathcal H_K}^2=x_i^\top K_T^{-1}x_i$. Since $f_i$ itself is a feasible interpolation and
    $\|f_i\|_{\mathcal H_K}\le B$, the minimum-norm property gives
    $\|\hat f_i^{\mathrm{MN}}\|_{\mathcal H_K}
    \le
    \|f_i\|_{\mathcal H_K}
    \le B$. Thus, $\hat f_i^{\mathrm{MN}}$ is also feasible for the constrained
    least-squares problem and achieves its minimum loss, namely zero.
    Consequently, in the noiseless setting, we may take
    $\mathcal F_{\mathrm{FTP}}(x_i)
    =
    \hat f_i^{\mathrm{MN}}$. This observation allows the individual fitting step in the FTP procedure to be represented in closed form through the kernel Gram matrix.
\end{remark}